% =============================================================================
\section{Results}
\label{sec:results}
% =============================================================================

We organize our results around three research questions:
\begin{itemize}[leftmargin=*]
\item \textbf{RQ1}: Does bottom-up (worst) selection achieve different effectiveness than top-down (best) selection?
\item \textbf{RQ2}: Does dual-end selection improve effectiveness, and at what cost?
\item \textbf{RQ3}: Does bidirectional ensemble improve over individual methods?
\end{itemize}

\subsection{Main Results}

Tables~\ref{tab:main_flan}, \ref{tab:main_qwen3}, and \ref{tab:main_qwen35} present the main results across all configurations, models, and datasets.

% Main results table with all 9 models
\begin{table*}[t]
\caption{Main results for the Flan-T5 family on TREC DL 2019 and 2020. Best values are in \textbf{bold}; second-best values are \underline{underlined}. All settings use \texttt{num\_child}=3, $k=10$, hits=100, and PL=128.}
\label{tab:main_flan}
\scriptsize
\setlength{\tabcolsep}{3pt}
\renewcommand{\arraystretch}{0.9}
\resizebox{\textwidth}{!}{%
\begin{tabular}{l|cccc|cccc|cccc}
\toprule
& \multicolumn{4}{c|}{\textbf{Flan-T5-Large (780M)}} & \multicolumn{4}{c|}{\textbf{Flan-T5-XL (3B)}} & \multicolumn{4}{c}{\textbf{Flan-T5-XXL (11B)}} \\
& \multicolumn{2}{c}{DL 2019} & \multicolumn{2}{c|}{DL 2020} & \multicolumn{2}{c}{DL 2019} & \multicolumn{2}{c|}{DL 2020} & \multicolumn{2}{c}{DL 2019} & \multicolumn{2}{c}{DL 2020} \\
\textbf{Method} & N@10 & MAP & N@10 & MAP & N@10 & MAP & N@10 & MAP & N@10 & MAP & N@10 & MAP \\
\midrule
\topdown{}-Heap & .6541 & .3074 & .6100 & .3511 & \underline{.6901} & .3172 & .6680 & .3696 & .6846 & .3201 & .6790 & .3727 \\
\topdown{}-Bubble & \textbf{.6874} & .3087 & \underline{.6264} & .3482 & \textbf{.6980} & \underline{.3182} & \textbf{.6868} & .3773 & \underline{.7077} & .3114 & \textbf{.6959} & .3693 \\
\bottomup{}-Heap & .2888 & .1737 & .2203 & .1786 & .6630 & .3060 & .6241 & .3511 & .6874 & .3211 & .6625 & .3625 \\
\bottomup{}-Bubble & .4571 & .2374 & .4116 & .2433 & .6730 & .2995 & .6691 & .3638 & .6936 & .3114 & .6811 & .3681 \\
\dualend{}-Cocktail & \underline{.6708} & \textbf{.3136} & \textbf{.6308} & \textbf{.3553} & .6884 & .3095 & \underline{.6795} & \textbf{.3798} & \textbf{.7137} & .3155 & \underline{.6895} & .3718 \\
\dualend{}-Selection & .6420 & .3016 & .6081 & .3350 & .6792 & .3057 & .6600 & .3750 & .6974 & .3155 & .6754 & .3787 \\
\bidir{}-RRF & .5820 & .2855 & .5580 & .3251 & .6845 & .3215 & .6509 & .3703 & .6905 & \textbf{.3265} & .6673 & .3726 \\
\bidir{}-Weighted & .6147 & .3062 & .5805 & .3371 & .6810 & .3176 & .6528 & .3673 & .6734 & .3215 & .6649 & .3694 \\
\bottomrule
\end{tabular}}
\end{table*}

\textbf{High-level pattern.} The DualEnd family is the strongest single-method family overall: \dualend{}-Cocktail wins 11 of 18 configurations, \dualend{}-Selection wins another 3 (Qwen3-4B on both DL19 and DL20 plus Qwen3-8B DL19), and \topdown{}-Bubble wins the remaining 4. \topdown{}-Heap, \bottomup{}-Heap, \bottomup{}-Bubble, \bidir{}-RRF, and \bidir{}-Weighted never win the headline NDCG@10 cell at the 4-decimal level. Mean NDCG@10 across the 9 models is \dualend{}-Cocktail $\geq$ \dualend{}-Selection $\geq$ \topdown{}-Bubble $\geq$ \topdown{}-Heap on both DL19 and DL20, with all four families well separated from \bottomup{} and \bidir{}.

\subsection{RQ1: Bottom-Up vs.\ Top-Down Effectiveness}
\label{sec:rq1}

We examine whether LLMs achieve different ranking effectiveness when asked to identify the \emph{least} relevant document compared to the \emph{most} relevant.

\textbf{Overall effectiveness}: Tables~\ref{tab:main_flan}, \ref{tab:main_qwen3}, and \ref{tab:main_qwen35} reveal a consistent pattern: \bottomup{} methods are substantially weaker than their \topdown{} counterparts across all models. On DL 2019 with Flan-T5-XL, \topdown{}-Heap achieves .6901 NDCG@10 while \bottomup{}-Heap achieves .6630---a gap of 2.7 points. On DL 2020, the gap widens to 4.4 points (.6680 vs.\ .6241).

The gap is dramatically \emph{larger} for smaller models: with Flan-T5-Large, \bottomup{}-Heap drops to .2888 on DL 2019 (vs.\ .6541 for \topdown{})---a catastrophic 36.5-point gap. For larger models, the gap narrows: with Qwen3-14B, \bottomup{}-Heap achieves .6966 (vs.\ .7447 for \topdown{}, a 4.8-point gap) and with Qwen3.5-27B, the gap is 3.1 points (.7135 vs.\ .7449). This scale-dependence suggests that worst-selection requires more model capacity than best-selection.

\textbf{Bubblesort vs.\ heapsort}: Within the \bottomup{} family, bubblesort outperforms heapsort for some larger models: on DL 2019 with Qwen3.5-27B, \bottomup{}-Bubble achieves .7336 versus .7135 for \bottomup{}-Heap. However, for smaller models the pattern reverses (T5-Large: .4571 vs.\ .2888). Bubblesort's $n-1$ forward passes provide more opportunities for error correction through repeated comparisons of neighboring elements, but at $\sim$1683 comparisons per query it is more than twenty times the cost of \topdown{}-Heap.

\textbf{Significance.} Paired approximate randomization with Bonferroni correction over the 18 configurations confirms that \bottomup{}'s loss against the best TopDown baseline is statistically robust. The mean delta is $-0.0616$, the family is positive in 0/18 configurations, and 6 configurations exceed the corrected significance threshold: Flan-T5-Large DL19 ($-0.2302$, $p_\text{Bonf}<0.001$), Flan-T5-Large DL20 ($-0.2148$, $p_\text{Bonf}<0.001$), Qwen3-14B DL20 ($-0.0649$, $p_\text{Bonf}=0.001$), Qwen3.5-4B DL19 ($-0.0950$, $p_\text{Bonf}=0.005$), Qwen3.5-4B DL20 ($-0.0749$, $p_\text{Bonf}=0.016$), and Qwen3.5-9B DL19 ($-0.0570$, $p_\text{Bonf}=0.025$). No \bottomup{} variant produces a Bonferroni-significant win on any TREC DL configuration.

\textbf{Refuting the irrelevance-detection hypothesis}: Our initial hypothesis---that LLMs may be better at identifying clearly irrelevant documents---is decisively refuted on these 9 models. Worst-selection is harder than best-selection across all model sizes, and the difficulty gap is largest for the smallest models. We conjecture that ``least relevant'' requires the model to hold a coherent relevance ordering of all candidates and select the minimum, while ``most relevant'' often has a salient winner that stands out. The position-bias analysis in \S\ref{sec:bias} provides additional evidence: standalone worst-selection collapses heavily onto the last passage (40--63\% of the time across models), suggesting that the model defaults to a positional heuristic rather than a relevance judgment.

\begin{table*}[t]
\caption{Main results for the Qwen3 family on TREC DL 2019 and 2020. Best values are in \textbf{bold}; second-best values are \underline{underlined}. All settings use \texttt{num\_child}=3, $k=10$, hits=100, and PL=512.}
\label{tab:main_qwen3}
\scriptsize
\setlength{\tabcolsep}{3pt}
\renewcommand{\arraystretch}{0.9}
\resizebox{\textwidth}{!}{%
\begin{tabular}{l|cccc|cccc|cccc}
\toprule
& \multicolumn{4}{c|}{\textbf{Qwen3-4B (4B)}} & \multicolumn{4}{c|}{\textbf{Qwen3-8B (8B)}} & \multicolumn{4}{c}{\textbf{Qwen3-14B (14B)}} \\
& \multicolumn{2}{c}{DL 2019} & \multicolumn{2}{c|}{DL 2020} & \multicolumn{2}{c}{DL 2019} & \multicolumn{2}{c|}{DL 2020} & \multicolumn{2}{c}{DL 2019} & \multicolumn{2}{c}{DL 2020} \\
\textbf{Method} & N@10 & MAP & N@10 & MAP & N@10 & MAP & N@10 & MAP & N@10 & MAP & N@10 & MAP \\
\midrule
\topdown{}-Heap & .6775 & .3115 & \underline{.6488} & .3479 & .6819 & .3241 & .6532 & .3591 & .7447 & .3420 & .6962 & .3783 \\
\topdown{}-Bubble & .6491 & .2904 & .6269 & .3288 & .6794 & .3034 & .6372 & .3433 & .7455 & .3354 & \underline{.7044} & .3875 \\
\bottomup{}-Heap & .6261 & .2916 & .5963 & .3236 & .6431 & .2897 & .5963 & .3262 & .6966 & .3244 & .6323 & .3463 \\
\bottomup{}-Bubble & .6305 & .2765 & .5778 & .3060 & .6273 & .2857 & .5948 & .3123 & .6702 & .3027 & .6395 & .3491 \\
\dualend{}-Cocktail & .6796 & .3064 & .6454 & .3436 & \underline{.7155} & .3239 & \textbf{.6678} & .3653 & \textbf{.7519} & .3403 & \textbf{.7051} & \textbf{.3907} \\
\dualend{}-Selection & \textbf{.7220} & \textbf{.3265} & \textbf{.6627} & \textbf{.3484} & \textbf{.7158} & .3256 & .6583 & .3607 & \underline{.7475} & \textbf{.3457} & .7003 & .3863 \\
\bidir{}-RRF & \underline{.6814} & .3369 & .6245 & .3504 & .6826 & .3225 & \underline{.6600} & \textbf{.3765} & .7172 & .3431 & .6736 & .3775 \\
\bidir{}-Weighted & .6608 & .3072 & .6376 & .3479 & .6784 & .3197 & .6397 & .3537 & .7200 & .3423 & .6741 & .3746 \\
\bottomrule
\end{tabular}}
\end{table*}


\subsection{RQ2: Dual-End Effectiveness and Efficiency}
\label{sec:rq2}

We assess whether \dualend{} can improve effectiveness, and at what cost.

\textbf{Parsing success rate}: Reliable dual output parsing is critical. For Flan-T5 models, we use likelihood scoring internally for the dual-end comparison---reading both the maximum and minimum of the softmax distribution over passage labels in a single forward pass. This avoids output parsing entirely and achieves 100\% success rate. For Qwen3/3.5 models, a cascading parser handles varied output formats including bracketed labels (\texttt{[A]}), numeric labels, and verbose responses. The overall parse failure rate (requiring heuristic defaults) is below 1\% across all models. The worst case is Qwen3-14B on DL 2020 with a 0.9\% partial parse rate (254 out of ${\sim}$29K comparisons).

\textbf{Effectiveness}: The \dualend{} family is the strongest overall. \dualend{}-Cocktail ranks first in mean NDCG@10 across all 9 models on both DL 2019 and DL 2020, and it is the strongest \emph{single} method in 11 of 18 model--dataset configurations across Tables~\ref{tab:main_flan}, \ref{tab:main_qwen3}, and \ref{tab:main_qwen35}. \dualend{}-Selection wins in 3 additional configurations (Qwen3-4B DL19, Qwen3-4B DL20, Qwen3-8B DL19), so DualEnd methods jointly achieve the highest NDCG@10 in 14 of 18 configurations. On DL 2019, representative gains include .7137 with T5-XXL (vs.\ .7077 for \topdown{}-Bubble), .7519 with Qwen3-14B (vs.\ .7455), and .7475 with Qwen3.5-27B (vs.\ .7449). The gain over the best top-down baseline ranges from $+0.0026$ to $+0.0446$ NDCG@10 points depending on model and dataset, with one exception (\dualend{}-Cocktail vs.\ \topdown{}-Bubble on Flan-T5-Large DL19 is $-0.0165$).

\dualend{}-Selection is especially competitive on smaller Qwen models, achieving .7220 with Qwen3-4B on DL 2019 and .6627 on DL 2020, where it outperforms all other methods by margins of $+0.0445$ and $+0.0139$ respectively. More broadly, the selection sort variant tends to favor smaller Qwen models, while cocktail shaker sort is more robust across model scales and families.

\textbf{Significance.} Across the 18 TREC DL configurations, the DualEnd family is positive against the best TopDown baseline in 14 cases (mean delta $+0.0058$). However, only one configuration survives Bonferroni correction: Qwen3-4B DL 2019, where \dualend{}-Selection beats \topdown{}-Heap by $+0.0446$ ($p_\text{Bonf}=0.010$, 95\% CI $[+0.0216, +0.0695]$, 28-13-2 wins-losses-ties). A near-miss is Qwen3-8B DL 2019 with $+0.0340$ ($p_\text{Bonf}=0.727$, 26-12-5). The remaining 16 configurations are consistent with the directional pattern but their bootstrap intervals cross zero, so we interpret DualEnd as the strongest empirical family overall while not claiming a uniform statistically significant improvement on these 43/54-query test sets.

\textbf{Why does DualEnd work?} The cocktail shaker sort alternates between forward passes (sinking the worst) and backward passes (bubbling the best). Each comparison identifies both extremes, so each pass makes progress at \emph{both} ends of the ranking. This dual progress is the key advantage: while \topdown{}-Bubble only moves the best document per pass, \dualend{}-Cocktail simultaneously moves the worst to the bottom---effectively getting two sorting operations per LLM call. The position-bias analysis in \S\ref{sec:bias} adds a second mechanism: the dual prompt restructures the model's positional defaults, with \dualend{}'s ``best'' selection becoming more uniform and its ``worst'' selection developing a reversed primacy pattern relative to standalone \bottomup{}.

\textbf{Efficiency}: The efficiency story of \dualend{} is nuanced. Table~\ref{tab:efficiency} compares the methods.

\begin{table}[!htbp]
\caption{Efficiency comparison for top-$k{=}10$ ranking of $n{=}100$ documents with $c{=}3$. Theoretical complexity and empirical measurements (per query, averaged across DL 2019).}
\label{tab:efficiency}
\small
\resizebox{\columnwidth}{!}{%
\begin{tabular}{lcccc}
\toprule
\textbf{Method} & \textbf{Theory} & \textbf{Comps} & \textbf{Tokens} & \textbf{Time} \\
\midrule
\topdown{}-Heap & $O(n{+}k\log_c n)$ & 77 & 32K & 7s \\
\topdown{}-Bubble & $O(kn/c)$ & 314 & 129K & 27s \\
\bottomup{}-Heap & $O(n \log_c n)$ & 273 & 112K & 23s \\
\bottomup{}-Bubble & $O((n{-}1)n/c)$ & 1683 & 691K & 146s \\
\dualend{}-Cocktail & $O(kn/c)$ & 546 & 224K & 30s \\
\dualend{}-Selection & $O(\frac{k}{2} \cdot \frac{n}{c{+}1})$ & 406 & 159K & 27s \\
\bidir{}-RRF & $2 \times$ Heap & 350 & 144K & 30s \\
\bottomrule
\end{tabular}
}
\end{table}

\dualend{}-Cocktail uses 546 comparisons---about $7\times$ \topdown{}-Heap's 77 and $1.7\times$ \topdown{}-Bubble's 314---but achieves higher mean NDCG@10. The additional comparisons are not wasted: each dual comparison produces \emph{two} ranking decisions, and the cocktail shaker sort exploits both to make progress from both ends of the ranking. \dualend{}-Selection uses 406 comparisons (29\% fewer than cocktail) with competitive effectiveness, especially on the smallest Qwen models.

The token-level overhead is similar: \dualend{}-Cocktail uses 224K total tokens per query versus 129K for \topdown{}-Bubble, a 74\% increase. We therefore frame \dualend{} explicitly as a quality-first refinement rather than a free improvement: the headline contribution is ``more bits per call extracted in a structurally useful way,'' not ``fewer total calls.'' The Pareto analysis in \S\ref{sec:pareto} formalizes this and identifies a missing intermediate point between \topdown{}-Bubble and \dualend{}-Cocktail that motivates the refinement variants in \S\ref{sec:refinement}.

\begin{table*}[t]
\caption{Main results for the Qwen3.5 family on TREC DL 2019 and 2020. Best values are in \textbf{bold}; second-best values are \underline{underlined}. All settings use \texttt{num\_child}=3, $k=10$, hits=100, and PL=512.}
\label{tab:main_qwen35}
\scriptsize
\setlength{\tabcolsep}{3pt}
\renewcommand{\arraystretch}{0.9}
\resizebox{\textwidth}{!}{%
\begin{tabular}{l|cccc|cccc|cccc}
\toprule
& \multicolumn{4}{c|}{\textbf{Qwen3.5-4B (4B)}} & \multicolumn{4}{c|}{\textbf{Qwen3.5-9B (9B)}} & \multicolumn{4}{c}{\textbf{Qwen3.5-27B (27B)}} \\
& \multicolumn{2}{c}{DL 2019} & \multicolumn{2}{c|}{DL 2020} & \multicolumn{2}{c}{DL 2019} & \multicolumn{2}{c|}{DL 2020} & \multicolumn{2}{c}{DL 2019} & \multicolumn{2}{c}{DL 2020} \\
\textbf{Method} & N@10 & MAP & N@10 & MAP & N@10 & MAP & N@10 & MAP & N@10 & MAP & N@10 & MAP \\
\midrule
\topdown{}-Heap & .7087 & .3289 & .6521 & .3678 & .7329 & \underline{.3418} & \underline{.6950} & .3870 & \underline{.7449} & \textbf{.3521} & .7104 & .4010 \\
\topdown{}-Bubble & \underline{.7108} & .3258 & \underline{.6713} & .3755 & \underline{.7349} & \textbf{.3415} & .6925 & .3847 & .7435 & .3481 & \underline{.7178} & .4036 \\
\bottomup{}-Heap & .6158 & .2881 & .5550 & .3076 & .6779 & .3194 & .6319 & .3470 & .7135 & .3357 & .6426 & .3786 \\
\bottomup{}-Bubble & .6120 & .2850 & .5963 & .3263 & .6712 & .3051 & .6677 & .3700 & .7336 & .3378 & .7004 & .3971 \\
\dualend{}-Cocktail & \textbf{.7161} & .3210 & \textbf{.6768} & \textbf{.3773} & \textbf{.7370} & .3356 & \textbf{.6984} & .3920 & \textbf{.7475} & .3492 & \textbf{.7186} & \textbf{.4075} \\
\dualend{}-Selection & .7022 & \textbf{.3288} & .6487 & .3471 & .7309 & .3399 & .6927 & \textbf{.3942} & .7319 & .3456 & .7122 & .4036 \\
\bidir{}-RRF & .6614 & .3167 & .6403 & .3547 & .7101 & .3347 & .6647 & .3722 & .7198 & .3457 & .6781 & .3980 \\
\bidir{}-Weighted & .6714 & .3204 & .6567 & .3638 & .7087 & .3349 & .6768 & .3798 & .7229 & .3458 & .6871 & .3995 \\
\bottomrule
\end{tabular}}
\end{table*}

\FloatBarrier

\subsection{RQ3: Bidirectional Ensemble}
\label{sec:rq3}

We evaluate whether fusing top-down and bottom-up rankings improves over either individual method.

\textbf{Ensemble effectiveness}: Contrary to our hypothesis, \bidir{} does \emph{not} outperform the \topdown{} baseline. On DL 2019 with Flan-T5-XL, \bidir{}-RRF achieves .6845 and \bidir{}-Weighted achieves .6810---both \emph{below} \topdown{}-Heap (.6901). This pattern holds across all 9 models and both datasets. The best BiDir result on DL 2019 is .7229 with Qwen3.5-27B, still below \topdown{}-Heap's .7449. Per-query analysis reveals that fusion ``beats the best individual'' on only 8--11 queries while it lands ``worse than the worst individual'' on 10--11 queries---a net negative effect.

\textbf{Significance.} The mean BiDir delta against the best TopDown baseline is $-0.0232$. The family is positive in 3/18 configurations (Qwen3-4B DL19 +0.0039, Qwen3-8B DL19 +0.0008, Qwen3-8B DL20 +0.0068), but none of these are Bonferroni-significant. By contrast, BiDir produces 3 Bonferroni-significant losses: Flan-T5-Large DL19 ($-0.0727$, $p_\text{Bonf}=0.013$), Flan-T5-Large DL20 ($-0.0459$, $p_\text{Bonf}=0.035$), and Qwen3.5-27B DL20 ($-0.0307$, $p_\text{Bonf}=0.016$). The asymmetry between (no significant wins) and (multiple significant losses) is the strongest evidence that direction-based ensembling does not work in this regime.

\textbf{Why does BiDir fail to improve?} The weakness of the \bottomup{} component drags down the ensemble. Since \bottomup{}-Heap is consistently 2.7--36.5 points below \topdown{}-Heap, fusing these two rankings dilutes the stronger signal. Even with $\alpha=0.7$ (giving 70\% weight to the top-down component), the bottom-up noise cannot be fully suppressed. Ranking agreement analysis (\S\ref{sec:agreement}) confirms: TopDown and DualEnd share Kendall's $\tau \approx 0.93$ (high agreement), while TopDown and BottomUp only reach $\tau \approx 0.86$---the divergence is not complementary, it is noise.

\textbf{Fusion weight sensitivity}: We test $\alpha \in \{0.3, 0.5, 0.7, 0.9\}$ (Table~\ref{tab:ablation_alpha}). Consistently across all models, $\alpha=0.9$ (90\% TopDown, 10\% BottomUp) is optimal. The monotonic improvement with increasing $\alpha$ confirms that \bottomup{} contributes only noise; the limit $\alpha\to 1$ is just plain TopDown.

\textbf{Comparison with permutation voting}: We compare \bidir{} (2$\times$ cost) with permutation voting ($p=2$, also 2$\times$ cost). Permutation voting uses the same 2$\times$ budget but applies it to the same algorithm with different orderings, addressing position bias rather than directional complementarity. On Qwen3-14B DL 2019, \bidir{}-RRF (.7172) loses to \topdown{}-Heap (.7447), while DualEnd-Cocktail beats \bidir{}-RRF on 23 of 43 queries at lower cost. This reinforces the conclusion that, when extra compute is available, redirecting it into joint best-worst elicitation (DualEnd) is more useful than redirecting it into a worst-only ranking pass (BiDir).

\textbf{Negative result significance}: This negative result challenges the intuitive assumption that combining rankings from different ``perspectives'' will always help. The key difference from permutation voting~\citep{zhuang2024setwise} is that permutation voting combines rankings from the \emph{same} algorithm with different input orders (addressing position bias), while \bidir{} combines rankings from \emph{different} algorithms where one is fundamentally weaker. Directional complementarity is not analogous to order-based complementarity.


\subsection{Ablation Studies}
\label{sec:ablation}

\textbf{Effect of window size} (\texttt{num\_child}): Table~\ref{tab:ablation_nc} shows the effect of varying $c \in \{2, 3, 5, 7\}$ on \dualend{}-Cocktail effectiveness on DL 2019.

\begin{table}[!htbp]
\caption{Effect of \texttt{num\_child} ($c$) on \dualend{}-Cocktail (DL 2019). Optimal $c$ depends on model family.}
\label{tab:ablation_nc}
\small
\resizebox{\columnwidth}{!}{%
\begin{tabular}{cccccc}
\toprule
$c$ & Window & T5-XL & Qwen3-8B & Q3.5-9B & \#Comps \\
\midrule
2 & 3-way & .6988 & .7187 & .7392 & 815 \\
3 & 4-way & .6884 & .7155 & .7370 & 546 \\
5 & 6-way & .6749 & .7224 & .7336 & 333 \\
7 & 8-way & .6480 & .7249 & .7386 & 247 \\
\bottomrule
\end{tabular}
}
\end{table}

The effect of window size is model-family-dependent. For Flan-T5-XL (512-token context limit), smaller windows are better: $c=2$ (.6988) significantly outperforms $c=7$ (.6480), a 5.1-point gap. Larger windows cause more prompt truncation, degrading both best and worst selection accuracy. For Qwen3-8B (32K+ context), larger windows are slightly better: $c=7$ (.7249) edges $c=2$ (.7187). For Qwen3.5-9B, performance is stable across window sizes (.7336--.7392). Importantly, larger windows reduce comparison count (247 at $c=7$ vs.\ 815 at $c=2$), offering a quality-efficiency trade-off.

\textbf{Effect of passage length}: Table~\ref{tab:ablation_pl} reports the effect of passage truncation length on both \topdown{}-Heap and \dualend{}-Cocktail for one T5 model and two Qwen models.

\begin{table}[!htbp]
\caption{Effect of passage truncation length on NDCG@10 (DL 2019). T5-XL has a 512-token context limit; Qwen3-8B and Qwen3.5-9B have 32K+.}
\label{tab:ablation_pl}
\small
\resizebox{\columnwidth}{!}{%
\begin{tabular}{ccccccc}
\toprule
PL & \multicolumn{2}{c}{T5-XL} & \multicolumn{2}{c}{Qwen3-8B} & \multicolumn{2}{c}{Q3.5-9B} \\
 & TD-H & DE-C & TD-H & DE-C & TD-H & DE-C \\
\midrule
64 & .6821 & .6846 & .6772 & .7035 & .7207 & .7142 \\
128 & .6901 & .6884 & .6934 & .7168 & .7378 & .7406 \\
256 & .6890 & .6918 & .6950 & .7130 & .7329 & .7393 \\
512 & .6890 & .6879 & .6895 & .7122 & .7329 & .7370 \\
\bottomrule
\end{tabular}
}
\end{table}

For Flan-T5-XL, performance plateaus at PL=128 (the original paper's default), with PL=256 and PL=512 producing nearly identical results---the model's 512-token context limit means longer passages are truncated. For Qwen models, PL=128 provides a small boost over PL=64 (+1.3--1.7 points), but PL=256 and PL=512 offer negligible further improvement. Crucially, the \emph{relative} advantage of \dualend{}-Cocktail over \topdown{}-Heap is consistent across all passage lengths, confirming our findings are not artifacts of a particular truncation setting.

\textbf{Fusion weight sensitivity} (\bidir{}-Weighted): For \bidir{}-Weighted, we vary $\alpha \in \{0.3, 0.5, 0.7, 0.9\}$ to control the relative importance of top-down versus bottom-up. Across all three representative models, $\alpha=0.9$ (90\% TopDown weight) is consistently optimal, confirming that the \bottomup{} component adds noise rather than complementary signal.

\begin{table}[!htbp]
\caption{Effect of fusion weight $\alpha$ for \bidir{}-Weighted on DL 2019 (NDCG@10). $\alpha{=}1.0$ is pure \topdown{}; $\alpha{=}0.0$ is pure \bottomup{}.}
\label{tab:ablation_alpha}
\small
\begin{tabular}{ccccc}
\toprule
$\alpha$ & 0.3 & 0.5 & 0.7 & 0.9 \\
\midrule
T5-XL & .6814 & .6845 & .6810 & .7002 \\
Q3-8B & .6618 & .6638 & .6784 & .6898 \\
Q3.5-9B & .6984 & .7087 & .7087 & .7240 \\
\bottomrule
\end{tabular}
\end{table}
